\documentclass[11pt]{article}

\usepackage[margin=0.85in]{geometry}
\usepackage{booktabs}
\usepackage{amsmath}
\usepackage{xcolor}

\definecolor{darkgreen}{RGB}{0,120,0}
\definecolor{darkred}{RGB}{180,0,0}
\newcommand{\worked}[1]{\textcolor{darkgreen}{\textbf{#1}}}
\newcommand{\failed}[1]{\textcolor{darkred}{\textbf{#1}}}

\setlength{\parskip}{0.3em}

% Tighter section spacing without pulling in titlesec (not in this TeX install).
\makeatletter
\renewcommand\section{\@startsection{section}{1}{\z@}%
  {-1.6ex \@plus -0.4ex \@minus -0.2ex}{0.7ex \@plus 0.2ex}%
  {\normalfont\large\bfseries}}
\renewcommand\subsection{\@startsection{subsection}{2}{\z@}%
  {-1.2ex \@plus -0.3ex \@minus -0.2ex}{0.5ex \@plus 0.2ex}%
  {\normalfont\normalsize\bfseries}}
\makeatother

\title{\vspace{-2em}Uncertainty Signals for Handwritten Math: Initial Results\\
\large Qwen2.5-VL-3B on FERMAT}
\author{}
\date{}

\begin{document}
\maketitle
\vspace{-2em}

\section*{The question}

Can a model tell us when it is wrong, without an answer key? We test one idea:
sample the model's answer $K=5$ times and measure how much it disagrees with
itself (Shannon entropy over the clustered answers). If disagreement tracks
error, it is a free error detector at deployment time. Two tasks:
\textbf{transcribing} a handwritten math solution, and \textbf{grading} whether
that solution contains a mistake.

\section*{Setup}

Qwen2.5-VL-3B on FERMAT. $K=5$ samples at temperature 0.7 per prompt, plus 2
greedy samples as a determinism check. Entropy is computed only from the
model's own outputs; correctness is scored separately against ground truth.
AUROC measures whether the two line up --- concretely, how often a randomly
chosen wrong item has higher entropy than a randomly chosen correct one.
Intervals are 95\% bootstrap CIs (10{,}000 resamples). $n=300$ items, sampled
to a balanced 150/150 split on whether the solution contains an error, with
thresholds registered before the run.

\section*{Results}

\begin{center}
\begin{tabular}{lccl}
\toprule
Task & AUROC (95\% CI) & Model accuracy & Verdict \\
\midrule
Transcription & \textbf{0.835} $[0.787, 0.879]$ & 0.393 & \worked{entropy predicts error} \\
Grading & 0.522 $[0.462, 0.582]$ & 0.517 (chance = 0.500) & \failed{no signal} \\
\bottomrule
\end{tabular}
\end{center}

\subsection*{Transcription: the signal is real \worked{(works)}}

When the model cannot read the handwriting it genuinely wavers, and that
wavering tracks its errors. The result survives every check applied: removing
all parse failures and all maximum-entropy items simultaneously still leaves
AUROC $=0.796$ $[0.736, 0.852]$ at $n=206$. It replicates an earlier $n=100$
pilot (0.788) on a like-for-like comparison.

One directly usable rule falls out of it: \textbf{when all five transcriptions
disagree, the model is wrong 57 times out of 61 (93.4\% precision)}, catching
31\% of all transcription errors. That is an abstention rule requiring no
ground truth. Two real items show what the signal looks like:

\begin{center}
\small
\begin{tabular}{@{}p{0.16\linewidth}p{0.78\linewidth}@{}}
\toprule
\multicolumn{2}{@{}l}{\textbf{Entropy 0 --- five identical readings, and correct}} \\
\midrule
ground truth & \texttt{option d} \\
5 samples & \texttt{option d}, \texttt{option d}, \texttt{option d}, \texttt{option d}, \texttt{option d} \\
\midrule
\multicolumn{2}{@{}l}{\textbf{Maximum entropy --- five different readings, and wrong}} \\
\midrule
ground truth & $A - B \neq B - A$ \\
5 samples & $A - B = \{1,3,5\}$ and $B - A = \{8\}$ \quad|\quad $A = B \neq B - A$ \quad|\quad
             $A = B \neq B = A$ \quad|\quad $A \neq B$ \quad|\quad $A \neq B \neq B - A$ \\
\bottomrule
\end{tabular}
\end{center}

\noindent The second item is genuinely hard handwriting; the model cannot settle
on a reading and every attempt differs. That instability is the signal.

Notably, only 3 of 300 items were both fully consistent \emph{and} scored wrong
--- and inspecting all three, each is a scoring artifact rather than a model
error (e.g.\ the model answering \texttt{option a} against a ground truth
written as \texttt{\textbackslash text\{A\}}). So the model is almost never
confidently wrong here, and the 0.393 accuracy above is an undercount.

\subsection*{Grading: no signal, because the model cannot do the task \failed{(fails)}}

The reason matters more than the number. FERMAT is 84.8\% error-containing.
On a sample inheriting that imbalance the model scored 75\% --- which looks
reasonable until you notice that always answering ``there is an error,''
without looking at anything, scores 87\%. The model was \emph{below} the
do-nothing baseline. On a balanced set the flattery disappears:

\begin{center}
\begin{tabular}{lc}
\toprule
Grading accuracy (balanced) & \textbf{0.517} \quad (chance $=0.500$) \\
\quad on solutions that do contain an error & 0.947 \\
\quad on solutions that are correct & \textbf{0.087} \\
Answers ``there is an error'' & 93\% of all items \\
\bottomrule
\end{tabular}
\end{center}

The model has one firm opinion --- that everything contains a mistake --- and
repeats it, so it is perfectly self-consistent almost everywhere: 42\% of items
had all five samples identical, and only three distinct entropy values occurred
across the entire run. There is nothing for entropy to measure.

\section*{What this means}

Entropy measures \emph{self-consistency}, never correctness. The method only
works when the two align --- when a model is steady while right and wavers while
wrong. That holds for transcription here. It fails for grading, because a
confidently wrong model is indistinguishable from a confidently right one: both
look steady.

The clearest evidence is a sub-result: on the correct solutions, grading AUROC
is \textbf{0.239} --- below 0.5, meaning entropy ran \emph{backwards}. The model
was at its most consistent exactly when it was wrong, insisting on errors that
were not there. Low entropy there is actively misleading, not merely uninformative.

So the headline is not ``entropy does not work for reasoning.'' It is
\textbf{sampling-based uncertainty inherits the model's competence}: it can flag
errors for a model that knows something, and tells you nothing about a model
that is guessing. Establishing that required rebalancing the evaluation set;
on the original imbalanced data the problem was invisible.

\section*{Next step}

The grading conclusion is about \emph{this model}, not the method. The natural
test is to repeat the grading arm with a stronger model (7B, or a text-only
model given the ground-truth transcription) --- gated on a capability check
first: measure grading accuracy against the 0.500 baseline \emph{before}
computing any uncertainty metric. If it still grades at chance, the negative
result stands for this task family; if it clears roughly 0.65--0.70, the
entropy question becomes answerable and the existing analysis applies unchanged.

The transcription result is deployable as-is and deserves an operating-point
analysis (precision/recall across thresholds, conformal calibration). That needs
no new model runs.

\vspace{0.5em}
\noindent\rule{\linewidth}{0.4pt}\\
{\small Full methodology, the two failed approaches that preceded the working
transcription metric, the $n=100$ pilot, and the diagnostic work behind the
grading result are in \texttt{report/report.pdf}. All analysis code is unit
tested in \texttt{pilot/}.}

\end{document}
